\documentclass[11pt,a4paper]{article}

\usepackage[utf8]{inputenc}
\usepackage[T1]{fontenc}
\usepackage[italian]{babel}
\usepackage{amsmath,amssymb}
\usepackage{geometry}
\usepackage{hyperref}

\geometry{margin=2.5cm}
\setlength{\parindent}{0pt}
\setlength{\parskip}{0.6em}
\hypersetup{
  colorlinks=true,
  linkcolor=black,
  urlcolor=blue
}

\title{Verifica e Correzioni dei Requisiti 1 e 2}
\author{Progetto Online Learning Applications}
\date{}

\begin{document}

\maketitle

\section{Sintesi}

Sono stati verificati i requisiti 1 e 2 del progetto rispetto al codice del
branch \texttt{giulia}. Il requisito 1 era gia' sostanzialmente soddisfatto. Il
requisito 2 era implementato solo parzialmente, perche' la versione precedente
del Combinatorial-UCB poteva generare azioni realizzate che violavano il conflict
graph. Il codice e' stato corretto introducendo una distribuzione su azioni
congiunte ammissibili.

\section{Requirement 1: singola campagna in ambiente stocastico}

Il requisito 1 richiede:
\begin{itemize}
  \item un ambiente stocastico per una singola campagna;
  \item una strategia UCB1 che ignora il vincolo di budget;
  \item una strategia UCB-like che gestisce il vincolo di budget.
\end{itemize}

\subsection{Stato del codice}

Il codice soddisfa il requisito. La classe \texttt{SingleCampaignEnv}
implementa una first-price auction in cui il massimo bid concorrente e' generato
come massimo di bid i.i.d. uniformi. Di conseguenza, la probabilita' di vincita
per un bid fissato e' descritta dalla distribuzione Beta usata nel clairvoyant.

La classe \texttt{UCB1BiddingAgent} implementa UCB1 ignorando i costi. La classe
\texttt{UCBLikeBiddingAgent} mantiene stime ottimistiche della reward, stime
pessimistiche del costo e risolve un LP per scegliere una distribuzione sui bid
ammissibile rispetto al budget medio
\[
  \rho = \frac{B}{T}.
\]

\subsection{Miglioramenti consigliati}

Per rendere la valutazione piu' completa, e' utile riportare separatamente un
caso con budget non vincolante e uno con budget vincolante, mostrando sia regret
sia costo cumulato. Questa scelta e' gia' presente nella pipeline
\texttt{run\_req1.py}.

\section{Requirement 2: campagne multiple in ambiente stocastico}

Il requisito 2 richiede:
\begin{itemize}
  \item un ambiente stocastico multi-campagna;
  \item una distribuzione congiunta sui massimi bid concorrenti;
  \item una strategia Combinatorial-UCB con vincolo di budget.
\end{itemize}

\subsection{Problema individuato}

La prima versione del codice risolveva un LP con variabili marginali
\(x_{ik}\), una per ogni coppia campagna--bid, e aggiungeva vincoli per budget e
conflitti. Tuttavia la scelta finale veniva campionata indipendentemente per
ogni campagna. Questo poteva produrre azioni realizzate non ammissibili:
campagne collegate da un arco di conflitto potevano ricevere entrambe un bid
nello stesso round.

Il problema era ancora piu' evidente nella fase di fallback iniziale, dove
l'agente sceglieva greedy una bid per ogni campagna, ignorando i conflitti.

\subsection{Correzione applicata}

La classe \texttt{CombinatorialUCBAgent} ora enumera lo spazio delle azioni
congiunte ammissibili. Ogni azione congiunta specifica, per ciascuna campagna,
un bid oppure l'astensione \(-1\), e viene costruita in modo da rispettare
sempre il conflict graph.

L'LP non sceglie piu' marginali indipendenti, ma una distribuzione \(\gamma\) su
azioni congiunte:
\[
  \max_{\gamma} \sum_{a \in \mathcal{A}} \gamma_a \, \widehat{f}^{UCB}(a)
\]
soggetta a
\[
  \sum_{a \in \mathcal{A}} \gamma_a \, \widehat{c}^{LCB}(a) \leq \rho,
  \qquad
  \sum_{a \in \mathcal{A}} \gamma_a = 1,
  \qquad
  \gamma_a \geq 0.
\]

Poiche' ogni \(a \in \mathcal{A}\) e' gia' fattibile, i vincoli di
incompatibilita' sono rispettati in ogni round realizzato, non solo in valore
atteso.

\subsection{Correzione del clairvoyant}

Anche il clairvoyant multi-campagna e' stato aggiornato per risolvere un LP
sullo stesso spazio di azioni congiunte ammissibili. In questo modo il regret
empirico confronta l'agente con un benchmark coerente con i vincoli del
problema.

\subsection{Correzione dell'ambiente}

L'ambiente \texttt{MultiCampaignEnv} ora rifiuta esplicitamente una bid vector
che assegna bid simultaneamente a due campagne incompatibili. Questo rende
visibile qualsiasi violazione del modello invece di correggerla a posteriori
scartando una delle due vittorie.

\section{Conclusione}

Dopo la correzione, il requisito 2 e' soddisfatto in modo piu' fedele:
\begin{itemize}
  \item la distribuzione dei massimi bid concorrenti e' una distribuzione
  congiunta prodotto sulle campagne;
  \item il budget e' gestito tramite LP;
  \item le campagne incompatibili non possono piu' ricevere bid nello stesso
  round;
  \item il clairvoyant e l'agente usano lo stesso spazio di azioni ammissibili.
\end{itemize}

\end{document}
